% ============================================================
\section{Advanced and Fallback Power-Flow Solvers}
\label{sec:advanced-pf}
% ============================================================

This chapter documents the remaining power-flow solver families declared in
\texttt{include/hacdcpf/power\_flow/}: the Continuation Power Flow
(\texttt{CpfSolver}), the Homotopy Continuation Solver, the Distributed Slack
Solver mathematical model, and the Newton--Krylov GMRES fallback.  It also
provides the complete \texttt{RobustNonlinearOptions} and
\texttt{PowerFlowOptions} reference tables.

% ------------------------------------------------------------
\subsection{Continuation Power Flow (CPF) / Voltage Stability}
\label{sec:cpf}
% ------------------------------------------------------------

The \texttt{CpfSolver} traces the P--V curve of a parameterized power system
from the base operating point (\(\lambda=0\)) toward the voltage-collapse nose
point.  The implementation lives in
\texttt{src/power\_flow/homotopy\_continuation.cpp} (the CPF corrector) and
\texttt{include/hacdcpf/power\_flow/voltage\_stability.hpp}.  It is exposed at
the lower-level \texttt{hacdcpf::powerflow} namespace and is not part of the
public \texttt{hacdcpf.hpp} facade.

\subsubsection{Parameterized Power System}

The loading direction is described by \texttt{CpfDirection}:
\begin{align}
P_{L,i}(\lambda) &= P_{L,i}^0 + \lambda\,\Delta P_{L,i}, \\
Q_{L,i}(\lambda) &= Q_{L,i}^0 + \lambda\,\Delta Q_{L,i}, \\
P_{G,i}(\lambda) &= P_{G,i}^0 + \lambda\,\Delta P_{G,i},\quad \text{(non-slack only)}.
\end{align}
Here \(\lambda=0\) is the base case; \(\lambda=1\) doubles all loads when
using the proportional direction.  Two factory methods build standard directions:
\begin{itemize}
  \item \texttt{CpfDirection::proportional(data)} --- all buses scale in
  proportion to their base-case power.
  \item \texttt{CpfDirection::at\_buses(data, bus\_indices)} --- uniform
  1~MW/Mvar per bus in the specified list.
\end{itemize}

\subsubsection{Predictor--Corrector Continuation}

The predictor step uses a secant (or constant) approximation:
\begin{align}
x_{\lambda+\Delta\lambda} \approx x_\lambda + \Delta\lambda\,\frac{x_\lambda - x_{\lambda-\Delta\lambda}}{\Delta\lambda_{\mathrm{prev}}}.
\end{align}

The corrector step solves the parameterized Newton system at the predicted point
by running the hybrid Newton PF with the modified injection vector and
convergence tolerance \texttt{corrector\_tol}.

Step adaptation:
\begin{align}
\Delta\lambda_{\mathrm{new}} =
\begin{cases}
\min(\texttt{step\_grow}\cdot\Delta\lambda,\;\texttt{step\_max}), & \text{corrector converged,} \\
\texttt{step\_shrink}\cdot\Delta\lambda, & \text{corrector failed.}
\end{cases}
\end{align}
If \(\Delta\lambda < \texttt{step\_min}\), the solver reports the nose point
and stops.

\subsubsection{\texttt{CpfOptions} Reference}

\begin{longtable}{L{0.26\linewidth}L{0.13\linewidth}L{0.48\linewidth}}
\toprule
Field & Default & Meaning \\
\midrule
\endhead
\texttt{lambda\_max} & 5.0 & Maximum loading factor \(\lambda\) to trace. \\
\texttt{step\_init} & 0.05 & Initial continuation step \(\Delta\lambda\). \\
\texttt{step\_min} & 1e-5 & Minimum step size; stop and declare nose point reached. \\
\texttt{step\_max} & 0.20 & Maximum step size (prevents skipping over the nose). \\
\texttt{step\_grow} & 2.0 & Growth factor applied after a successful corrector solve. \\
\texttt{step\_shrink} & 0.5 & Shrink factor applied after a corrector failure. \\
\texttt{corrector\_max\_iter} & 50 & Max Newton iterations in each corrector call. \\
\texttt{corrector\_tol} & 1e-8 & Convergence tolerance for the corrector (p.u.). \\
\texttt{monitor\_bus} & -1 & Bus index to plot on the P--V curve; $-1$ = auto-select max demand. \\
\texttt{trace\_all\_buses} & false & Store full \texttt{vm}/\texttt{va} vectors at every CPF step. \\
\texttt{max\_steps} & 2000 & Safety cap on total predictor--corrector steps. \\
\texttt{vm\_min\_pu} & 0.25 & Voltage floor; stop if any bus falls below this (p.u.). \\
\texttt{use\_secant\_predictor} & true & Use secant (two-point) predictor; false = constant predictor. \\
\bottomrule
\end{longtable}

\subsubsection{CPF Result Structure}

Each tracing step produces a \texttt{CpfPoint}:
\begin{itemize}
  \item \texttt{lambda} -- continuation parameter value at this step;
  \item \texttt{vm\_monitor} -- monitor-bus voltage magnitude (p.u.).
\end{itemize}
The top-level result includes the full trace of \texttt{CpfPoint} records, the
final converged solution, estimated nose-point \(\lambda^*\) and minimum
voltage, and the convergence status.

\subsubsection{Algorithm A13: CPF Predictor--Corrector (condensed)}
\label{algo:a13-inline}
The full pseudo-code and flowchart appear in Section~\ref{algo:a13}.
The condensed form below highlights the inner step logic.

\begin{algorithmic}[1]
\Procedure{CpfSolver::solve}{$\mathit{base}$, $\mathit{dir}$, $\mathit{opts}$}
  \State Solve base case at $\lambda{=}0$ via \texttt{NewtonSolver} (A3)
  \State $\lambda \gets 0$;\quad $\Delta\lambda \gets \mathit{opts.step\_init}$
  \While{$\lambda < \lambda_{\max}$ and iter $<$ max\_steps}
    \State Predictor: secant extrapolation $\to x_{\mathrm{pred}}$ at $\lambda + \Delta\lambda$
    \State Corrector: Newton from $x_{\mathrm{pred}}$ at $\lambda_{\mathrm{next}}$
    \If{converged}
      \State Accept; $\lambda \gets \lambda_{\mathrm{next}}$;\quad
        $\Delta\lambda \gets \min(\mathit{step\_grow}\cdot\Delta\lambda,\,\mathit{step\_max})$
    \Else
      \State $\Delta\lambda \gets \mathit{step\_shrink}\cdot\Delta\lambda$
      \If{$\Delta\lambda < \mathit{step\_min}$} declare nose; \textbf{break} \EndIf
    \EndIf
  \EndWhile
  \State \Return \textit{CpfResult}\{trace, nose~$\lambda^*$, status\}
\EndProcedure
\end{algorithmic}

% ------------------------------------------------------------
\subsection{Homotopy Continuation Solver}
\label{sec:homotopy}
% ------------------------------------------------------------

The \texttt{HomotopyContinuationSolver} solves the hybrid power flow by
gradually ramping ``difficult'' injections from a flat-start base to the target
values.  It is invoked as a last-resort fallback when the direct Newton solve
fails (Phase~4 of the robust pipeline).

\subsubsection{Homotopy Path}

The homotopy path scales all active/reactive loads and converter setpoints:
\begin{align}
P_{L}(\lambda) &= \lambda\,P_{L}^{\mathrm{target}}, &
Q_{L}(\lambda) &= \lambda\,Q_{L}^{\mathrm{target}}, \\
P_{\mathrm{vsc}}(\lambda) &= \lambda\,P_{\mathrm{vsc}}^{\mathrm{target}}, &
P_{\mathrm{dcdc}}(\lambda) &= \lambda\,P_{\mathrm{dcdc}}^{\mathrm{target}}.
\end{align}
At \(\lambda=0\) the system is trivially solved from the flat start.
At \(\lambda=1\) the full target system is solved.

\subsubsection{Step Adaptation}

\begin{align}
\text{Success:}\quad\Delta\lambda &\leftarrow \min(2\,\Delta\lambda,\,\Delta\lambda_{\max}), \\
\text{Failure:}\quad\Delta\lambda &\leftarrow \frac{\Delta\lambda}{2};\quad
\text{fail if}\;\Delta\lambda < \Delta\lambda_{\min}.
\end{align}
The current implementation records accepted/rejected steps in a
\texttt{HomotopyState} diagnostic struct and returns both the final
\texttt{PowerFlowResult} and the state for instrumentation.

\subsubsection{Algorithm A14: Homotopy Continuation (condensed)}
\label{algo:a14-inline}
See Section~\ref{algo:a14} for the full pseudo-code and flowchart.

\begin{algorithmic}[1]
\Procedure{HomotopyContinuationSolver::solve}{$\mathit{sys}$, $\mathit{opt}$}
  \State Scale to $\lambda{=}0$; solve trivially
  \While{$\lambda < 1$}
    \State $\lambda_{\mathrm{next}} \gets \min(1,\,\lambda + \Delta\lambda)$;\quad
      scale trial system
    \If{Newton (A3) from current solution converges}
      \State Accept; $\lambda \gets \lambda_{\mathrm{next}}$;\quad $\Delta\lambda \gets \min(2\Delta\lambda,\,\Delta\lambda_{\max})$
    \Else
      \State $\Delta\lambda \gets \Delta\lambda/2$
      \If{$\Delta\lambda < \Delta\lambda_{\min}$} \Return failed \EndIf
    \EndIf
  \EndWhile
  \State \Return PowerFlowResult at $\lambda{=}1$
\EndProcedure
\end{algorithmic}

% ------------------------------------------------------------
\subsection{Distributed Slack Solver}
\label{sec:distributed-slack}
% ------------------------------------------------------------

The distributed-slack solver distributes the active-power mismatch across a set
of participating buses according to participation factors instead of assigning
it all to a single slack bus.  Two variants are exposed:
\begin{itemize}
  \item \texttt{solve\_power\_flow\_distributed\_slack} -- simplified form
  (one Newton solve followed by a reported allocation);
  \item \texttt{solve\_power\_flow\_distributed\_slack\_full} -- bounded
  redispatch with repeated complete Newton solves.  The legacy ``full'' name is
  retained for API compatibility; this is not an augmented-Jacobian method.
\end{itemize}

\subsubsection{Problem Formulation}

Let $\alpha_g \ge 0$ be participation factors (normalized to sum to 1) and let
$P_s$ be the total system mismatch.  Each participating generator $g$ receives
an additional dispatch increment:
\begin{align}
\delta P_g = \alpha_g P_s,
\end{align}
subject to individual upper bounds \(\texttt{max\_participation\_p}_g\) when
provided.  The reference bus ($\texttt{reference\_bus}$) provides the angle
reference; it need not be the same as the swing generator.

\subsubsection{Convergence Criterion}

The simplified variant inherits convergence from its single Newton solve and
reports the resulting mismatch allocation without changing the solved state.
The bounded variant repeatedly solves Newton PF, redispatches the remaining
active mismatch subject to generator and participation limits, and stops when
the unallocated correction is at most \(\max(10^{-10},\varepsilon)\) p.u.  It
fails closed when the limits cannot absorb the mismatch or after 20 solves.

\subsubsection{Result Structure}

\texttt{DistributedSlackResult} extends the standard bus-voltage vectors with:
\begin{itemize}
  \item \texttt{distributed\_slack\_p} -- map from participating bus index to
  final active correction (p.u. on the system \texttt{base\_mva});
  \item \texttt{hit\_limits} -- list of buses that reached
  \texttt{max\_participation\_p} and were clamped.
\end{itemize}

\subsubsection{Algorithm A7: Distributed-Slack PF (condensed)}
\label{algo:a7-inline}
See Section~\ref{algo:a7} for the full pseudo-code and flowchart.

\begin{algorithmic}[1]
\Procedure{solve\_distributed\_slack}{$\mathit{sys}$, $\mathit{cfg}$, $\mathit{opt}$}
  \State Apply \texttt{reference\_bus} as the angle reference
  \Repeat
    \State Run Newton PF; compute the unified-injection mismatch $\Delta P_s$
    \State Allocate $\Delta P_s$ by normalized $\alpha_g$ and clamp to generator limits
    \State Redispatch participating generators
  \Until{$|\Delta P_s|\le\max(10^{-10},\varepsilon)$ p.u. or 20 solves}
  \State \Return \textit{DistributedSlackResult}
\EndProcedure
\end{algorithmic}

% ------------------------------------------------------------
\subsection{Newton--Krylov GMRES Fallback (Phase 5)}
\label{sec:newton-krylov}
% ------------------------------------------------------------

Phase~5 of the robust pipeline replaces the direct-factorization linear solve
with a GMRES iterative solver inside the Newton loop.  The implementation lives
in \texttt{src/power\_flow/newton\_krylov.cpp} and is declared in
\texttt{include/hacdcpf/power\_flow/newton\_krylov.hpp}.  It is opt-in and not
active by default (\texttt{RobustNonlinearOptions::enable\_krylov\_fallback}).

At each Newton iteration the augmented linear system
\begin{align}
(J + \alpha M^{-1})\,\Delta x = -F
\end{align}
is solved via GMRES with a Schur-complement AC/DC block preconditioner $M$:
\begin{align}
M =
\begin{bmatrix}
J_{ac} & 0 \\
0 & J_{dc}
\end{bmatrix},
\end{align}
where $J_{ac}$ is the AC subsystem Jacobian and $J_{dc}$ is the DC subsystem
Jacobian.  The Krylov fallback is most effective when the Jacobian is ill-
conditioned due to strong AC/DC coupling, but adds per-iteration overhead
compared to the direct SparseLU path.

\subsubsection{Algorithm A15: Newton--Krylov GMRES (condensed)}
\label{algo:a15-inline}
See Section~\ref{algo:a15} for the full pseudo-code and flowchart.

\begin{algorithmic}[1]
\Procedure{NewtonKrylovSolver::solve}{$\mathit{data}$, $\mathit{opt}$}
  \State Compute $F(x)$; form $J$; build Schur preconditioner $M = \mathrm{blkdiag}(J_{\mathrm{AC}}, J_{\mathrm{DC}})$
  \While{$\|F\|_\infty \ge \varepsilon$ and iter $<$ max}
    \State GMRES($m$): solve $(J + \alpha M^{-1})\Delta x = -F$ iteratively
    \State $x \gets x + \Delta x$;\quad re-evaluate $F(x)$
  \EndWhile
  \State \Return \textit{NewtonKrylovResult}
\EndProcedure
\end{algorithmic}

% ------------------------------------------------------------
\subsection{Complete \texttt{RobustNonlinearOptions} Reference}
\label{sec:robust-options-table}
% ------------------------------------------------------------

\begin{longtable}{L{0.35\linewidth}L{0.12\linewidth}L{0.42\linewidth}}
\toprule
Field & Default & Description \\
\midrule
\endhead
\multicolumn{3}{l}{\textbf{Phase~1: Scaling and diagnostics}}\\
\midrule
\texttt{enable\_residual\_scaling} & \texttt{true} & Scale residual equations by \(\max|F_i^{\mathrm{spec}}|\). \\
\texttt{enable\_variable\_scaling} & \texttt{true} & Scale Newton variables (\(\theta,V,V^{dc}\)) to \(O(1)\). \\
\texttt{enable\_jacobian\_row\_col\_equilibration} & \texttt{true} & Apply \(D_F J D_x^{-1}\) row/column equilibration. \\
\texttt{enable\_condition\_monitor} & \texttt{true} & Compute \(\ell_1\) row/col proxy for condition number each iteration. \\
\texttt{min\_scale} & 1e-6 & Floor for any scale factor. \\
\texttt{max\_scale} & 1e6 & Ceiling for any scale factor. \\
\texttt{min\_vm\_pu} & 1e-8 & Voltage magnitude floor (pu) in Jacobian \(\partial P/\partial V\). \\
\texttt{min\_branch\_x\_pu} & 1e-20 & Min branch reactance; below this, branch excluded from FDPF \(B'\). \\
\texttt{bad\_condition\_threshold} & 1e10 & Proxy condition number above which Jacobian is flagged ``bad''. \\
\texttt{tiny\_pivot\_threshold} & 1e-12 & Pivot magnitude threshold for informational logging. \\
\midrule
\multicolumn{3}{l}{\textbf{Phase~2: Nonmonotone line search}}\\
\midrule
\texttt{enable\_nonmonotone\_linesearch} & \texttt{true} & Use GLL nonmonotone Armijo search with sliding window. \\
\texttt{nonmonotone\_window} & 5 & Sliding window \(M\) for nonmonotone reference merit. \\
\texttt{armijo\_c} & 1e-4 & Armijo sufficient-decrease constant \(c_1 \in (0,1)\). \\
\texttt{line\_search\_beta} & 0.5 & Step-length reduction factor \(\beta\) per backtrack trial. \\
\texttt{max\_line\_search\_trials} & 12 & Maximum backtrack trials before step rejection. \\
\midrule
\multicolumn{3}{l}{\textbf{Phase~2: Smooth Fischer--Burmeister NCP (off by default)}}\\
\midrule
\texttt{enable\_smooth\_ncp} & \texttt{false} & Use \(\mu\)-perturbed smooth-FB for PV/PQ complementarity. \\
\texttt{ncp\_mu0} & 1e-2 & Initial smoothing parameter \(\mu_0\). \\
\texttt{ncp\_mu\_min} & 1e-12 & Minimum smoothing parameter (limit of continuation). \\
\texttt{ncp\_mu\_factor} & 0.1 & Fine-phase annealing factor (near-convergence). \\
\texttt{ncp\_mu\_factor\_coarse} & 0.5 & Coarse-phase annealing factor (large residual). \\
\texttt{ncp\_mu\_phase\_transition} & 0.1 & Fraction of initial residual to switch to fine annealing. \\
\texttt{enable\_activity\_hysteresis} & \texttt{true} & Hold PV/PQ active-set for \texttt{min\_active\_set\_hold\_iters} iterations. \\
\texttt{min\_active\_set\_hold\_iters} & 3 & Minimum iterations between consecutive PV/PQ status changes. \\
\texttt{q\_limit\_hysteresis} & 1e-4 & Hysteresis band on Q-limit proximity (pu) before switching. \\
\midrule
\multicolumn{3}{l}{\textbf{Phase~3: LM trust-region fallback}}\\
\midrule
\texttt{enable\_lm\_trust\_region\_fallback} & \texttt{true} & Compute LM step \((J^\top J+\lambda I)\Delta x=-J^\top F\) on Newton failure. \\
\texttt{lm\_lambda0} & 1e-4 & Initial LM damping parameter \(\lambda_0\). \\
\texttt{lm\_lambda\_min} & 1e-12 & Minimum LM damping. \\
\texttt{lm\_lambda\_max} & 1e8 & Maximum LM damping. \\
\texttt{trust\_region\_radius0} & 1.0 & Initial trust-region radius. \\
\texttt{trust\_region\_radius\_min} & 1e-6 & Minimum radius; below this the LM fallback is abandoned. \\
\texttt{trust\_region\_radius\_max} & 100.0 & Maximum trust-region radius. \\
\midrule
\multicolumn{3}{l}{\textbf{Phase~3: PTC-SER pseudo-transient continuation}}\\
\midrule
\texttt{enable\_ptc\_ser} & \texttt{true} & Use Self-Equalising Residual (SER) exponent for adaptive PTC. \\
\texttt{ptc\_dt0} & 0.1 & Initial pseudo-timestep \(\delta t_0\). \\
\texttt{ptc\_dt\_min} & 1e-4 & Minimum pseudo-timestep (large damping). \\
\texttt{ptc\_dt\_max} & 1e6 & Maximum pseudo-timestep (approaches pure Newton). \\
\midrule
\multicolumn{3}{l}{\textbf{Phase~4: Homotopy continuation (see Section~\ref{sec:homotopy})}}\\
\midrule
\texttt{enable\_homotopy\_fallback} & configurable & Trigger \texttt{HomotopyContinuationSolver} when Phase~3 fails. \\
\midrule
\multicolumn{3}{l}{\textbf{Phase~5: Krylov fallback (off by default)}}\\
\midrule
\texttt{enable\_krylov\_fallback} & \texttt{false} & Replace SparseLU with Newton--GMRES + Schur-complement preconditioner. \\
\bottomrule
\end{longtable}

% ------------------------------------------------------------
\subsection{\texttt{PowerFlowOptions} Complete Reference}
\label{sec:pf-options-table}
% ------------------------------------------------------------

The complete field-by-field reference for \texttt{PowerFlowOptions} and its
seven orthogonal sub-structs (\texttt{ConvergenceOptions},
\texttt{LinearSolverOptions}, \texttt{GlobalizationOptions},
\texttt{PVPQSwitchingOptions}, \texttt{ConverterModeOptions},
\texttt{ZipLoadOptions}, \texttt{RuntimeOptions}) is in
Section~\ref{sec:pf-options-sub}.  Table~\ref{tab:pf-options-flat} below
lists only the fields that are directly relevant to the advanced solver paths
described in this chapter.

\begin{longtable}{L{0.38\linewidth}L{0.14\linewidth}L{0.37\linewidth}}
\toprule
Field & Default & Description \\
\midrule
\endhead
\texttt{ac\_eval\_threads} & 0 & Threads for AC Jacobian evaluation; 0 = hardware auto-detect (\texttt{std::thread::hardware\_concurrency()}). \\
\texttt{globalization} & \texttt{LineSearch} & \texttt{LineSearch}, \texttt{TrustRegion}, or \texttt{PseudoTransient}. \\
\texttt{trust\_region\_radius0} & 1.0 & Initial trust-region radius (used when \texttt{TrustRegion} is selected). \\
\texttt{trust\_region\_max} & 10.0 & Maximum trust-region radius. \\
\texttt{ptc\_delta0} & 1.0 & Initial PTC pseudo-timestep. \\
\texttt{ptc\_growth} & 2.0 & PTC pseudo-timestep growth factor. \\
\texttt{enable\_semi\_smooth\_newton} & false & Use NCP-based PV/PQ switching (Fischer-Burmeister). \\
\texttt{enable\_coupled\_jacobian} & false & Include cross-coupling Jacobian blocks. \\
\texttt{enable\_augmented\_equations} & false & Use augmented-equation path for VDC/VAC bus setpoints. \\
\texttt{enable\_krylov\_fallback} & false & Replace SparseLU with Newton--GMRES + Schur-complement preconditioner. \\
\texttt{enable\_homotopy\_fallback} & configurable & Trigger \texttt{HomotopyContinuationSolver} when Phase~3 fails. \\
\bottomrule
\end{longtable}
\label{tab:pf-options-flat}
